\section{Efficiency bound and augmentation covariance}
\label{sec:supp-efficiency}

Fix \(M\) and \(h\). The observed-data law factors as
\[
p_O(o)
=
p_W(w)
\{\pi_M(w)p_{Y\mid W}(y\mid w)\}^{r}
\{1-\pi_M(w)\}^{1-r},
\]
where the factor involving \(Y\) is present only when \(r=1\). With known \(\pi_M\), the nonparametric tangent space is the \(L_2(P)\)-closure of
\begin{equation}
s(O)=s_W(W)+R s_Y(Y,W),
\qquad
E s_W(W)=0,
\qquad
E\{s_Y(Y,W)\mid W\}=0.
\label{eq:known-pi-tangent}
\end{equation}
When the verification law is unrestricted, the tangent space additionally contains
\begin{equation}
s_\pi(O)=\{R-\pi_M(W)\}c(W)
\label{eq:pi-tangent}
\end{equation}
for square-integrable \(c\). This is the observed-data tangent decomposition for missing-at-random models; see \citet{RobinsRotnitzkyZhao1994}, Chapters 3--4 of \citet{Tsiatis2006}, and Chapter 25 of \citet{VanderVaart1998}.

For a finite signed measure \(\vartheta\) on \(\mathcal Y\), let
\[
\zeta_\vartheta(Y)=\int_{\mathcal Y}1(Y\le y)\,d\vartheta(y),
\qquad
m_{0,\vartheta}(W)=E\{\zeta_\vartheta(Y)\mid W\},
\]
and
\[
\theta_{M,h,\vartheta}
=
\frac{E\{K_h(X-x_0)\zeta_\vartheta(Y)\}}{\mu_h}.
\]
Finite total variation of \(\vartheta\) makes \(\zeta_\vartheta\) bounded.

\begin{theorem}[Signed-projection canonical gradient]
\label{thm:signed-efficiency}
Under Condition~\ref{cond:sampling}, assume \(\mu_h>0\) and
\(E\{K_h^2(X-x_0)/\pi_M(W)\}<\infty\). The canonical gradient for
\(\theta_{M,h,\vartheta}\) is
\begin{equation}
\phi_{M,h,\vartheta}^{\mathrm{eff}}(O)
=
\frac{K_h(X-x_0)}{\mu_h}
\left[
 m_{0,\vartheta}(W)
 +\frac{R}{\pi_M(W)}\{\zeta_\vartheta(Y)-m_{0,\vartheta}(W)\}
 -\theta_{M,h,\vartheta}
\right].
\label{eq:supp-efficient-if}
\end{equation}
The gradient is orthogonal to the propensity-score tangent space \eqref{eq:pi-tangent}. For a finite threshold collection \(y_1,\ldots,y_L\) and \(\vartheta=\sum_{j=1}^Lc_j\delta_{y_j}\),
\[
\phi_{M,h,\vartheta}^{\mathrm{eff}}
=
\sum_{j=1}^Lc_j\phi_{M,h,y_j}^{\mathrm{eff}}.
\]
\end{theorem}

\begin{proof}
Consider a regular parametric submodel with score \(s=s_W+Rs_Y\) as in \eqref{eq:known-pi-tangent}. Write
\[
\mathcal N_{M,h,\vartheta}
=
E\{K_h(X-x_0)\zeta_\vartheta(Y)\},
\qquad
\theta_{M,h,\vartheta}=\mathcal N_{M,h,\vartheta}/\mu_h.
\]
Differentiating the numerator and denominator gives
\begin{align}
\dot{\mathcal N}_{M,h,\vartheta}
&=
E\{K_hm_{0,\vartheta}s_W\}
+
E(K_h\zeta_\vartheta s_Y),
\label{eq:num-path-derivative}\\
\dot\mu_h
&=E(K_hs_W).
\label{eq:den-path-derivative}
\end{align}
Hence
\begin{equation}
\dot\theta_{M,h,\vartheta}
=
\frac1{\mu_h}
\left[
E\{K_h(m_{0,\vartheta}-\theta_{M,h,\vartheta})s_W\}
+
E(K_h\zeta_\vartheta s_Y)
\right].
\label{eq:theta-path-derivative}
\end{equation}

The function in \eqref{eq:supp-efficient-if} has mean zero and lies in the closure of the tangent space. Its \(W\)-component is
\[
\frac{K_h}{\mu_h}(m_{0,\vartheta}-\theta_{M,h,\vartheta}),
\]
and its conditional-outcome component is
\[
\frac{K_h}{\mu_h\pi_M}(\zeta_\vartheta-m_{0,\vartheta}).
\]
For the \(W\)-score,
\[
E(\phi_{M,h,\vartheta}^{\mathrm{eff}}s_W)
=
\frac1{\mu_h}
E\{K_h(m_{0,\vartheta}-\theta_{M,h,\vartheta})s_W\}.
\]
Using \(R\perp Y\mid W\), \(E(s_Y\mid W)=0\), and \(R^2=R\),
\begin{align*}
E(\phi_{M,h,\vartheta}^{\mathrm{eff}}R s_Y)
&=
\frac1{\mu_h}
E\left[
K_h\frac{R}{\pi_M}(\zeta_\vartheta-m_{0,\vartheta})s_Y
\right]\\
&=
\frac1{\mu_h}E(K_h\zeta_\vartheta s_Y).
\end{align*}
The sum equals \eqref{eq:theta-path-derivative}, so \(\phi_{M,h,\vartheta}^{\mathrm{eff}}\) is the canonical gradient.

For a propensity score \(s_\pi\) in \eqref{eq:pi-tangent},
\begin{align*}
E\{\phi_{M,h,\vartheta}^{\mathrm{eff}}s_\pi\mid W\}
&=
\frac{K_hc(W)}{\mu_h}
E\left[
\frac{R}{\pi_M}(\zeta_\vartheta-m_{0,\vartheta})(R-\pi_M)
\,\Bigm|\,W
\right]\\
&=
\frac{K_hc(W)}{\mu_h}
\{1-\pi_M(W)\}
E(\zeta_\vartheta-m_{0,\vartheta}\mid W)
=0.
\end{align*}
Thus the canonical gradient is unchanged when \(\pi_M\) is included in the nuisance model. Linearity gives the finite-threshold vector result.
\end{proof}

Theorem~\ref{thm:efficiency} follows by taking \(\vartheta\) to be the coordinate point masses \(\delta_{y_1},\ldots,\delta_{y_L}\).

\begin{lemma}[Augmentation identity]
\label{lem:augmentation}
Under missing at random, for every integrable augmentation \(a_y(W)\),
\[
E\{U_y(a)\mid W\}=m_{0,y}(W),
\qquad
E\{U_y(a)\mid X\}=q_{0,y}(X).
\]
\end{lemma}

\begin{proof}
Conditional on \(W\),
\[
E(RZ_y\mid W)=\pi_M(W)m_{0,y}(W),
\qquad
E(R\mid W)=\pi_M(W).
\]
Therefore
\[
\begin{aligned}
E\{U_y(a)\mid W\}
&=a_y(W)+\frac1{\pi_M(W)}
E[R\{Z_y-a_y(W)\}\mid W]\\
&=a_y(W)+m_{0,y}(W)-a_y(W)
=m_{0,y}(W).
\end{aligned}
\]
Taking conditional expectation given \(X\) yields the second identity.
\end{proof}

\begin{lemma}[Conditional covariance decomposition]
\label{lem:cov-decomp}
Let \(\delta_y^a(W)=a_y(W)-m_{0,y}(W)\), and assume the displayed conditional second moments are finite. Then
\begin{equation}
\operatorname{Cov}\{U_y(a),U_z(a)\mid W\}
=
\frac{C_{yz}(W)}{\pi_M(W)}
+
\left\{\frac1{\pi_M(W)}-1\right\}
\delta_y^a(W)\delta_z^a(W).
\label{eq:supp-cov-decomp}
\end{equation}
\end{lemma}

\begin{proof}
Write
\begin{equation}
U_y(a)-m_{0,y}
=
\frac{R}{\pi_M}(Z_y-m_{0,y})
+
\delta_y^a\left(1-\frac{R}{\pi_M}\right).
\label{eq:decomp-proof}
\end{equation}
The two terms have conditional mean zero. Their cross moment vanishes because \(R\perp Y\mid W\) and \(E(Z_y-m_{0,y}\mid W)=0\). Since \(R^2=R\),
\[
E\left[
\frac{R^2}{\pi_M^2}
(Z_y-m_{0,y})(Z_z-m_{0,z})
\,\Bigm|\,W
\right]
=
\frac{C_{yz}(W)}{\pi_M(W)}.
\]
Also,
\[
E\left[\left(1-\frac{R}{\pi_M}\right)^2\Bigm|W\right]
=
\frac1{\pi_M(W)}-1.
\]
Substitution into \eqref{eq:decomp-proof} proves \eqref{eq:supp-cov-decomp}.
\end{proof}

\begin{proof}[Proof of Corollary~\ref{cor:cov-order}]
For a fixed augmentation \(a\), define
\[
\phi_{M,h,y}^{a}(O)
=
\frac{K_h}{\mu_h}
\{U_y(a)-F_{h,0}(y\mid x_0)\}.
\]
The difference from the canonical gradient is
\begin{equation}
\Delta_{M,h,y}^{a}(O)
=
\phi_{M,h,y}^{a}(O)-\phi_{M,h,y}^{m_0}(O)
=
\frac{K_h}{\mu_h}\delta_y^a(W)
\left(1-\frac{R}{\pi_M(W)}\right).
\label{eq:augmentation-orthogonal-component}
\end{equation}
For known \(\pi_M\), \(\Delta_{M,h,y}^{a}\) is orthogonal to the tangent space in \eqref{eq:known-pi-tangent}. For every \(s_W\),
\[
E(\Delta_{M,h,y}^{a}s_W\mid W)
=
\frac{K_h}{\mu_h}\delta_y^as_W
E\left(1-\frac{R}{\pi_M}\Bigm|W\right)=0.
\]
For every \(s_Y\) with \(E(s_Y\mid W)=0\),
\[
E(\Delta_{M,h,y}^{a}R s_Y\mid W)=0.
\]
Hence \(\phi_{M,h,y}^{a}\) is an influence function in the known-design model and \(\Delta_{M,h,y}^{a}\) is its orthogonal excess component.

When \(\pi_M\) is unrestricted,
\begin{equation}
E\{\Delta_{M,h,y}^{a}s_\pi\mid W\}
=
-\frac{K_h}{\mu_h}
\delta_y^a(W)c(W)\{1-\pi_M(W)\}.
\label{eq:augmentation-propensity-inner-product}
\end{equation}
This inner product vanishes for every \(c\) exactly when
\[
K_h(X-x_0)\{1-\pi_M(W)\}\delta_y^a(W)=0
\quad\text{a.s.}
\]
The oracle augmentation therefore retains the canonical-gradient property under both known and unknown verification laws, while the family \(\phi^a\) is the influence-function family for the known-design experiment.

By Lemma~\ref{lem:cov-decomp},
\[
\Omega_{M,h}^a(y,z)-\Omega_{M,h}^{m_0}(y,z)
=
\frac{h^d}{\mu_h^2}
E\left[
K_h^2
\left(\frac1{\pi_M}-1\right)
\delta_y^a\delta_z^a
\right].
\]
For thresholds \(y_1,\ldots,y_L\) and \(c\in\mathbb R^L\),
\[
\begin{aligned}
&\sum_{j=1}^L\sum_{\ell=1}^L
c_jc_\ell
\{\Omega_{M,h}^a(y_j,y_\ell)-\Omega_{M,h}^{m_0}(y_j,y_\ell)\}\\
&\quad=
\frac{h^d}{\mu_h^2}
E\left[
K_h^2
\left(\frac1{\pi_M}-1\right)
\left\{\sum_{j=1}^Lc_j\delta_{y_j}^a\right\}^2
\right]
\ge0.
\end{aligned}
\]
This proves the positive-semidefinite covariance decomposition.
\end{proof}

\begin{proof}[Proof of Theorem~\ref{thm:transfer}]
Along \(a_{\lambda,y}=q_y+\lambda D_y\),
\[
a_{\lambda,y}-m_{0,y}=-(A_y-\lambda D_y).
\]
Consequently,
\[
\begin{aligned}
J_{M,h}(\lambda)
&=E\{v_{M,h}\|A-\lambda D\|_\nu^2\}\\
&=E(v_{M,h}\|A\|_\nu^2)
-2\lambda E\{v_{M,h}\langle A,D\rangle_\nu\}
+\lambda^2E(v_{M,h}\|D\|_\nu^2)\\
&=J_{M,h}(0)-2\lambda G_{M,h}+\lambda^2H_{M,h}.
\end{aligned}
\]
If \(H_{M,h}=0\), weighted Cauchy--Schwarz gives
\[
|G_{M,h}|^2
\le
J_{M,h}(0)H_{M,h}=0.
\]
Thus \(G_{M,h}=0\), the criterion is constant on \([0,1]\), and the convention \(\lambda_h^{\mathrm{or}}=0\) applies.

If \(H_{M,h}>0\), the unconstrained minimizer is \(G_{M,h}/H_{M,h}\), and projection gives \(\lambda_h^{\mathrm{or}}\). Setting \(\lambda=1\) gives
\[
J_{M,h}(0)-J_{M,h}(1)=2G_{M,h}-H_{M,h}.
\]
The remaining cases in \eqref{eq:oracle-gain} follow according to whether the projected minimizer equals zero, lies in \((0,1)\), or equals one.
\end{proof}
